\documentclass[12pt]{article}
\usepackage{amsmath,amsfonts,amsthm}
\usepackage{fullpage}
\usepackage{algpseudocode}
\usepackage{algorithm}

\begin{document}

\author{Toby Hocking}

\title{Pool Adjacent Violator Algorithm (PAVA) for label position optimization}

\maketitle

\section{Introduction}

R packages like directlabels and animint2 implement aligned text label drawing in a way that uses optimization to ensure readability.
This paper describes the previous algorithm based on quadratic programming, and a new approach based on PAVA.

\section{Problem}

\subsection{Data}

We assume there are $n$ labels.
Each label $i\in\{1,\dots,n\}$ has two associated data: a target position $t_i\in\mathbb R$, and a half-size $h_i\in\mathbb R_+$.
We assume the labels are already sorted by target value, $t_1\leq\cdots\leq t_n$.

Let $x_1,\dots,x_n$ be the label centers (optimization variables).
For each label $i$, the lower limit of its text is $x_i-h_i$, and the upper limit is $x_i+h_i$.
We assume the user inputs a global lower bound $\underline B\in\mathbb R$ and upper bound $\overline B\in\mathbb R$ in which text can be drawn.

\subsection{Optimization}

The label positions are determined as the solution of the optimization problem:
\begin{equation}
    \label{eq:label-pos}
\begin{aligned}
  \min_{x\in\mathbb R^n} & ||x-t||^2 \\
  \text{subject to }& x_1 - \underline B\geq h_1 ,\\
                    & \forall i\in\{2, \dots, n\},\, x_i-x_{i-1} \geq h_i + h_{i-1},\\
   & \overline B-x_n\geq h_n .
\end{aligned}
\end{equation}

\section{Solutions}

\subsection{Previous solution using quadratic programming}
TODO.

\subsection{Proposed solution using PAVA}

The proposed solution involves using a change of variables and PAVA.
First, let the lower limit vector be
\begin{equation}
  \label{eq:L}
  L = \texttt{cumsum}(\underline B+h_1, h_1+h_2, \dots, h_{n-1}+h_n).
\end{equation}
Each element of the $L$ vector represents the lower feasible limit of the corresponding element of $x$, given the constraints.
In other words, $L\in\mathbb R^n$ is a vector, with elements defined as
\begin{equation}
  \label{eq:L_i}
  L_i =
  \begin{cases}
    \underline B+h_1 & \text{ for } i=1\\
    L_{i-1}+h_{i-1}+h_i & \text{ for }i\in\{2,\dots, n\}.
  \end{cases}
\end{equation}
The change of variables is targets $d=t-L$, and solution $v=x-L$.
In this new space, the feasible set is $v_i\in [0, U]$ for all $i$, where the upper limit is $U=\overline B-h_n-L_n$.

The corresponding optimization problem is
\begin{equation}
    \label{eq:pava-problem}
\begin{aligned}
  \min_{v\in\mathbb R^n} & ||v-d||^2 \\
  \text{subject to }& v_1 \geq 0 ,\\
                    & \forall i\in\{2, \dots, n\},\, v_{i-1} \leq v_{i},\\
   & v_n\leq U.
\end{aligned}
\end{equation}
By using the definitions of $d,v,L,U$, it can be shown that~(\ref{eq:pava-problem}) is equivalent to~(\ref{eq:label-pos}).

PAVA($d$, lower, upper) is an algorithm that computes a solution to~(\ref{eq:pava-problem}), given a vector of data $d$, and lower/upper bounds on feasible values for the solutions $v$.
After computing the PAVA solution $v$, the solution to~(\ref{eq:label-pos}) is computed via $x=v+L$.
Algorithm~\ref{alg:pava-pos} summarizes the solution method.

\begin{algorithm}
\begin{algorithmic}
  \Require half sizes $h_1,\dots,h_n\in\mathbb R$, target positions $t_1\leq\cdots\leq t_n\in\mathbb R$, bounds $\underline B<\overline B\in \mathbb R$.
  \State $L \gets \texttt{cumsum}(\underline B+h_1, h_1+h_2, \dots, h_{n-1}+h_n)$
  \State $U \gets \overline B-h_n-L_n$
  \State $v \gets \texttt{PAVA}(t-L, 0, U)$
  \State \Return $v+L$
\end{algorithmic}
  \caption{\label{alg:pava-pos}PAVA for label position optimization}
\end{algorithm}

\end{document}
